% Created (UTC): 2026-06-21T18:28:10Z
% Repository HEAD: 4d0a91f7c224c233697d6e5077ab6589475572e3
\documentclass[11pt]{article}
\usepackage[margin=1in]{geometry}
\usepackage{amsmath,amssymb,amsthm,mathtools}
\usepackage{booktabs}
\usepackage{microtype}
\emergencystretch=3em \hbadness=4000
\usepackage[colorlinks=true,linkcolor=blue,citecolor=blue,urlcolor=blue]{hyperref}
\DeclareMathOperator{\Li}{Li}
\DeclareMathOperator{\Cl}{Cl}
\newcommand{\iu}{\mathrm{i}}
\newcommand{\Imag}{\operatorname{Im}}
\newcommand{\Real}{\operatorname{Re}}
\theoremstyle{plain}\newtheorem{theorem}{Theorem}\newtheorem{proposition}[theorem]{Proposition}
\theoremstyle{remark}\newtheorem{remark}[theorem]{Remark}
\title{Clausen constants at roots of unity:\\
the fundamental basis by Dirichlet conductor, and a weight parity swap}
\author{Vladimir Reshetnikov}
\date{2026-06-21}
\begin{document}
\maketitle
\begin{center}\small
Research report. Created (UTC) 2026-06-21T18:28:10Z; repository \texttt{HEAD}
\texttt{4d0a91f7c224c233697d6e5077ab6589475572e3}.\\
Store \texttt{clausen-roots-of-unity.wl}. Companions: \texttt{polylog-vertical-line},
\texttt{cleo-arctan-sin-chi2}.
\end{center}

\begin{abstract}
The imaginary parts $\Cl_2(2\pi a/n)=\Imag\Li_2(e^{2\pi\iu a/n})$ that appear when one evaluates the
dilogarithm at algebraic points (as in Reshetnikov's Math.StackExchange questions 1430169 and 1410276)
are organized completely by the theory of odd Dirichlet characters. We record the explicit picture and
verify it numerically: the $\mathbb Q$-span of $\{\Cl_2(2\pi a/n)\}$ has dimension $\varphi(n)/2$, and the
genuinely \emph{new} constants entering at exact level $n$ are indexed by the \emph{primitive} odd
characters mod $n$. By conductor the fundamental constants are $V=\Cl_2(\pi/3)$ (Gieseking, conductor $3$),
$G=\Cl_2(\pi/2)$ (Catalan, conductor $4$), $W_8=\Cl_2(\pi/4)$ (conductor $8$), and $W_{24}=\Cl_2(\pi/12)$
(conductor $24$); conductor $12$ contributes nothing. This settles the open question in MSE answer
a/5113303: level $24$ needs exactly \emph{one} new constant $W_{24}$ beyond $\{G,V,W_8\}$. We give the
reduction tables, the $L$-value/trigamma forms, a new closed form
$\Li_2(e^{\iu\pi/12})=\tfrac{73\pi^2}{576}+\iu W_{24}$, and a clean \emph{weight parity swap}: at weight
$3$ the imaginary part $\Imag\Li_3(e^{2\pi\iu a/n})$ is \emph{elementary} (a rational multiple of $\pi^3$),
so the deep constants move to the real part---the reverse of weight $2$.
\end{abstract}

\section{Context: dilogarithms at algebraic points}
\href{https://math.stackexchange.com/q/1430169}{MSE 1430169} asks for a closed form of
$\Li_2\!\bigl(\tfrac12+\iu(1-\tfrac{\sqrt3}2)\bigr)$. The standard route (Sangchul Lee,
\href{https://math.stackexchange.com/a/1430262}{a/1430262}) uses $\Li_2(z)=-\Li_2(z/(z-1))-\tfrac12\ln^2(1-z)$
to send the point to the unit circle: here $z/(z-1)=e^{7\pi\iu/6}$, and
$\Li_2(e^{\iu\theta})=\bigl(\tfrac{\pi^2}6-\tfrac{\pi\theta}2+\tfrac{\theta^2}4\bigr)+\iu\Cl_2(\theta)$ for
$\theta\in[0,2\pi)$, so everything reduces to the Clausen value $\Cl_2(7\pi/6)$. The closed form's
nontrivial content is therefore the reduction of a single $\Cl_2$ at a rational angle. A later answer
(\href{https://math.stackexchange.com/a/5113303}{a/5113303}) expresses the whole level-$12$ family through
two constants---the Catalan constant $G$ and the Gieseking constant $V$---and asks whether level $24$
needs something new. It does; the organizing principle is classical, and we make it explicit.

\section{The fundamental constants, by conductor}
For $n>2$ the values $\Cl_2(2\pi a/n)$, $1\le a<n$, are linear combinations of $L(2,\chi)$ over the
\emph{odd} Dirichlet characters $\chi$ mod $n$ (those with $\chi(-1)=-1$), because $\Cl_2$ is an odd
function of the angle. Hence:
\begin{proposition}\label{prop:dim}
$\dim_{\mathbb Q}\operatorname{span}\{\Cl_2(2\pi a/n):1\le a<n\}=\varphi(n)/2$, the number of odd
characters mod $n$. The constants genuinely new at exact level $n$ (not already present at any proper
divisor) are indexed by the \emph{primitive} odd characters mod $n$; their count is
$0,0,1,0,2,1,2,0,4,0,5,0,1,2,7,\dots$ for $n=3,4,5,\dots$.
\end{proposition}
The first few fundamental constants, one per primitive odd character, are
\[
  \underbrace{V=\Cl_2(\tfrac\pi3)}_{\text{conductor }3},\quad
  \underbrace{G=\Cl_2(\tfrac\pi2)}_{\text{conductor }4},\quad
  \underbrace{W_8=\Cl_2(\tfrac\pi4)}_{\text{conductor }8},\quad
  \underbrace{W_{24}=\Cl_2(\tfrac\pi{12})}_{\text{conductor }24}.
\]
Conductor $12$ has \emph{no} primitive odd character (every odd character mod $12$ factors through $3$ or
$4$), which is exactly why the level-$12$ family of a/5113303 collapses onto $\{G,V\}$. Conductors $6$ and
$10$ likewise add nothing. The pentagon level $5$ is the first to add \emph{two} constants at once,
$\Cl_2(2\pi/5)$ and $\Cl_2(4\pi/5)$ (two conjugate primitive odd characters mod $5$).

Their canonical $L$-value (trigamma) forms are
\begin{equation}\label{eq:trig}
  G=\frac{\psi^{(1)}(\tfrac14)-\pi^2}{8},\qquad
  \sqrt3\,V=\frac{3\psi^{(1)}(\tfrac13)-2\pi^2}{6}\ \ \Bigl(\text{i.e. }\psi^{(1)}(\tfrac13)=\tfrac{2\pi^2}3+2\sqrt3\,V\Bigr).
\end{equation}

\section{Reduction tables and the settled level-24 question}
Every $\Cl_2(2\pi a/n)$ at $n\in\{8,12,24\}$ reduces to the basis above (PSLQ, $\ge110$ digits,
Champernowne canary; all confirmed in Wolfram):
\begin{align*}
\text{level }12:\quad
  &\Cl_2(\tfrac\pi6)=\tfrac23G+\tfrac14V,\quad \Cl_2(\tfrac{5\pi}6)=\tfrac23G-\tfrac14V,\quad
  \Cl_2(\tfrac{2\pi}3)=\tfrac23V;\\
\text{level }8:\quad
  &\Cl_2(\tfrac{3\pi}4)=W_8-\tfrac12G;\\
\text{level }24:\quad
  &\Cl_2(\tfrac{5\pi}{12})=\tfrac{16W_8+24W_{24}-4G-3V}{24},\quad
  \Cl_2(\tfrac{7\pi}{12})=\tfrac{4W_8+6W_{24}-3G}{6},\\
  &\Cl_2(\tfrac{11\pi}{12})=\tfrac{24W_{24}-8G-3V}{24}.
\end{align*}
The four level-$24$ values $\Cl_2(k\pi/12)$, $k\in\{1,5,7,11\}$, span a $4$-dimensional space
$\{G,V,W_8,W_{24}\}$ ($=\varphi(24)/2$), and $W_{24}$ is provably outside $\operatorname{span}\{G,V,W_8\}$.
\begin{remark}
This settles a/5113303: the dilogarithms of $24$th roots of unity require exactly one constant beyond the
level-$12$ pair $\{G,V\}$ and the level-$8$ constant $W_8$, namely $W_{24}=\Cl_2(\pi/12)$.
\end{remark}

\section{A new closed form}
Since $\Cl_2(\pi/12)=W_{24}$ is itself fundamental, the cleanest carrier is the dilogarithm at the
primitive $24$th root $e^{\iu\pi/12}$:
\begin{equation}\label{eq:new}
  \Li_2(e^{\iu\pi/12})=\frac{73\pi^2}{576}+\iu\,W_{24},
\end{equation}
the real part being $\tfrac{\pi^2}6-\tfrac{\pi}{2}\!\cdot\!\tfrac\pi{12}+\tfrac14(\tfrac\pi{12})^2$. Pulling
this back through $z=1/(1-e^{\iu\theta})$ gives Reshetnikov-style closed forms at the corresponding
$\mathbb Q(\sqrt2,\sqrt3,\sqrt6)$ points; the cleanest stays the root-of-unity form~\eqref{eq:new}.

\section{A weight parity swap}
The same character parity that makes $\Cl_2$ (weight $2$, odd characters) carry deep constants reverses one
weight up. For odd $\chi$, $L(s,\chi)$ is a rational multiple of $\pi^s$ exactly when $s$ is \emph{odd};
since $\Imag\Li_3(e^{\iu\theta})=\sum\sin(k\theta)/k^3$ pairs with odd characters at $s=3$,
\begin{equation}\label{eq:swap}
  \Imag\Li_3(e^{2\pi\iu a/n})\in\mathbb Q\cdot\pi^3\quad\text{for all }a,n,
\end{equation}
e.g. $\Imag\Li_3(e^{2\pi\iu/3})=\tfrac{2\pi^3}{81}$, $\Imag\Li_3(\iu)=\tfrac{\pi^3}{32}$,
$\Imag\Li_3(e^{\iu\pi/3})=\tfrac{5\pi^3}{162}$, $\Imag\Li_3(e^{2\pi\iu/5})=\tfrac{4\pi^3}{125}$. Thus at
weight $3$ the \emph{imaginary} part at roots of unity is elementary, and the deep constants
($\zeta(3)$ and the even-character $L(3,\chi)$) live in the \emph{real} part---the exact reverse of weight
$2$. (On the line $\Real z=1$ the picture is different again: there $\Imag\Li_3(1+\iu)$ is the central-binomial
Gaussian tower constant, because $1+\iu$ is not a root of unity; see \texttt{polylog-vertical-line}.)

\section{Verification and reproducibility}
Proposition~\ref{prop:dim}'s reductions, the trigamma forms~\eqref{eq:trig}, the new closed
form~\eqref{eq:new}, the original MSE value, and the weight-$3$ rationals~\eqref{eq:swap} are verified in
\texttt{mpmath} (PSLQ, $90$--$110$ digits, Champernowne canary) and Wolfram, residuals
$10^{-63}$ to $10^{-73}$. (Two level-$24$ reductions were sign-corrected against Wolfram before storing.)
Closed forms are in \texttt{src/PolyLog/identities/clausen-roots-of-unity.wl}; scripts under
\texttt{src/PolyLog/tools/scratch/clausen-roots/}.

\end{document}
